% 本文件由 LaTeX 排版 Agent 根据有效 translation.json 自动生成。
% author=Athanasius; work_id=2814; title=为逃亡辩护

\chapter{为逃亡辩护}

% structure: p0001 source=h1 target=omitted reason=page-title-already-rendered-by-parent

1. \Emph{\Person{亚大纳修}被指控因逃跑而怯懦。}\par

我听说，现今在\Place{安提约基}的\Person{莱翁提乌斯}、\Place{尼禄城}的\Person{那耳喀索斯}、现今在\Place{劳狄刻雅}的\Person{乔治}，以及和他们在一起的\OriginalTerm{亚略派}{Arian}人士，正在四处散播许多诽谤我的谣言，指控我怯懦，就因为当他们寻找我时，我没有自投其手。至于他们的指责和诽谤，尽管我可以写出许多连他们也无法否认的事，并且所有听说过他们所作所为的人都知道是真实的，但我不会被迫去回应他们，只愿提醒他们我主的话语，以及宗徒的宣言，即\BibleQuote{谎言是出于魔鬼}，以及\BibleQuote{辱骂人的，也不能承受天主的国}。因为这就足以证明，他们的思想和言语都不合乎福音，而是随从他们自己的喜好，凡他们自己所愿的，他们就以为是好的。\par

2. \Emph{这项指控的虚伪。}\par

但由于他们假装指控我怯懦，我有必要就此事写一些话，借此证明他们是心怀恶意的人，没有读过圣经；或者即使读过，也不相信其中所载神谕的天主默感。因为如果他们相信这一点，他们就不敢违背圣经行事，也不敢效法那些杀害了主的犹太人的恶意。因为天主曾赐给他们一条诫命：\BibleQuote{应孝敬你的父亲和你的母亲}，以及\BibleQuote{咒骂父亲或母亲的，应处以死刑}\BibleRef{玛 15:4}；那民族却另立了一条相反的律法，将孝敬变为不敬，并把子女本应奉献给父母的金钱，挪作他用。虽然他们读过\Person{达味}所行的事，他们的行为却与他的榜样相反，并且控告无辜的人，因为他们在安息日掐了麦穗，用手搓着吃。并非因为他们关心法律或安息日，因为他们在那天犯下了更严重地违背律法的罪；而是心怀恶意，嫉妒门徒得救的道路，并且希望自己的私意独尊。然而他们已获得了自己罪恶的报应，不再是一个圣洁的民族，从此被算作\Place{索多玛}的统治者，和\Place{哈摩辣}的百姓\BibleRef{依 1:10}。同样，这些人，在我看来，并不亚于他们，已因对自己的愚昧无知而受到了惩罚。因为他们不明白自己所说的是什么，却以为知道自己所不知道的事；而他们里面惟一的才智，就是行恶，并日复一日地策划更邪恶的计谋。因此，他们以我们目前的逃亡来指责我们，并非为了德行，好像希望我们挺身而出显示勇敢（敌人怎能为那些不与他们同奔狂途的人产生这种愿望呢？）；而是满怀恶意，假装如此，并四处散布这种说法，愚蠢地以为，我们因害怕他们的辱骂，就会被打动而自投其手。因为这正是他们所渴望的：为达此目的，他们用尽各种计谋：他们假装是朋友，却像敌人一样搜寻，为要使自己饱饮我们的血，并将我们铲除，因为我们始终反对并且仍在反对他们的不虔，驳斥并谴责他们的\OriginalTerm{异端}{heresy}。\par

3. \Emph{亚略派人士对主教们的暴行。}\par

他们迫害过谁，抓住过谁，而不随心所欲地侮辱伤害呢？他们搜寻过谁，找到过谁，而不使他要么悲惨地死去，要么遭受各种虐待呢？凡官长们似乎所做的，都是他们的作为；其他人只是他们意志和邪恶的工具。因此，有什么地方没有留下他们恶意的纪念呢？有谁曾反对他们，而他们不密谋陷害他，像\Person{依则贝耳}一样捏造借口毁掉他呢？有哪一个教会此刻不为他们针对主教们的阴谋得逞而哀哭呢？\Place{安提约基}为\OriginalTerm{正统的宣信者}{orthodox Confessor}\Person{欧斯塔提乌斯}哀悼；\Place{巴拉内阿}为最可钦佩的\Person{欧夫辣提翁}哀悼；\Place{帕尔图斯}和\Place{安塔拉杜斯}为\Person{基玛提乌斯}和\Person{卡尔特里乌斯}哀悼；\Place{哈德里安堡}为爱慕基督者\Person{欧特罗庇乌斯}及其继承人\Person{卢基乌斯}哀悼，后者常因他们的缘故身负锁链，就此丧亡；\Place{安基拉}为\Person{马尔切路斯}哀悼，\Place{贝罗亚}为\Person{基鲁斯}哀悼，\Place{加沙}为\Person{阿斯刻拉帕斯}哀悼。对于这些人，他们在施加许多暴行之后，又通过阴谋将他们放逐；但对\Place{色雷斯}的主教\Person{特奥杜洛}和\Person{奥林庇乌斯}，以及我们和我们的司铎，他们则竭力搜寻，意在一旦发现，就处以死刑；若非那时恰与他们意愿相反我们逃走了，恐怕我们早已如此丧命。因为授与此意的信件已经送达总督\Person{多纳图斯}，针对\Person{奥林庇乌斯}及其同人，也送达\Person{菲拉格里乌斯}，针对我。并且他们发动了对\Place{君士坦丁堡}主教\Person{保禄}的迫害，一找到他，就使人在\Place{卡帕多细亚}一个名叫\Place{库库苏斯}的地方将他公开绞死，为此目的雇用了总督\Person{菲利普}作为行刑者。他是他们异端的赞助者，也是他们邪恶计划的工具。\par

4. \Emph{米兰大公会议之后的行动。}\par

难道他们对这一切感到满意，并愿在未来保持安静吗？决不；他们尚未罢手，却如\BibleBook{}{箴言}中的水蛭，在邪恶中越发狂欢，且紧盯更大的教区。谁能充分描述他们已犯下的暴行？谁能数算他们所行的一切？就在最近，当教会处在和平中，人民在各自的集会中敬拜时，罗马主教\Person{利伯略}、高卢的都主教\Person{保利努斯}、意大利的都主教\Person{狄奥尼修斯}、撒丁岛群岛的都主教\Person{路西斐尔}，以及意大利的\Person{优西比乌斯}，所有这些好主教和真理的宣讲者，竟被逮捕并放逐，毫无任何借口，只因为他们不愿与\OriginalTerm{亚略异端}{Arian heresy}结合，也不肯签署他们捏造的针对我的虚假指控与诽谤。\par

5. \Emph{赞颂霍修斯}\par

关于那位名符其实的伟大\Person{霍修斯}，那位享有幸福老年的宣认者，我无需赘言，因为我想众人皆知，他们也使他遭受放逐；因为他并非无名之辈，而是所有人中最卓越的，而且不止于此。何时召开过公会议，他不曾居于领导地位，并以正确的劝导说服众人？哪座教会没有留下他牧职的光荣纪念？谁曾忧伤地来见他，而不曾欢欣离去？哪个穷乏者曾求他援助，而未得偿所愿？然而，就连这么一个人，他们也攻击，因为他知道他们为维护自己的不义而捏造的诽谤，故不愿赞同他们对我们的阴谋。纵然后来，在遭受了无以计数的鞭打以及针对他亲属的阴谋之后，他一度向他们屈服\BibleRef{迦 2:5}，因为他年老体弱，但即使如此，他们的邪恶也在这事上显露出来；他们如此竭力以一切方式证明自己并非真正的基督徒。\par

6. \Emph{乔治对亚历山大人的暴行}\par

此后，他们再次转向亚历山大，重新图谋杀害我们：他们的行径比以前更恶劣。因为教会突然间被士兵包围，祷告声被战斗声所取代。接着，他们派来的卡帕多细亚的\Person{乔治}，在四旬期抵达，带来了他们教他的更重的灾祸。因为在复活周过后，贞女们被投入监狱；主教们被士兵用锁链拖走；孤儿寡妇的家被掠夺，他们的食粮被拿走；住宅遭到袭击，基督徒在夜间被赶出，他们的住处被封；神职人员的兄弟因其弟兄而遭受生命危险。这些暴行已足够可怕，但更可怕的事随之而来。因为在圣神降临节后一周（5月11日），人民在斋戒后前往墓园祈祷，因为所有人都拒绝与乔治共融，那个无耻之人得知此事，便煽动指挥官\Person{塞巴斯蒂安}（他是个\OriginalTerm{摩尼教徒}{Manichee}）来对付他们；他立即率领大队士兵，手持武器、出鞘的剑、弓和矛，在主日攻击人民，尽管那天是主日：他发现只有少数人在祈祷（因为大部分人因时间已晚已经散去），便行出这等暴行，无愧于这些人的门徒。他点燃一堆火，将一些贞女置于火旁，试图强迫她们声明自己属于\OriginalTerm{亚略信仰}{Arian faith}；当他见她们不屈不挠，不惧烈火时，他立刻剥光她们的衣服，猛击她们的面部，致使她们几乎长时间无法辨认。\par

7. \Emph{乔治的暴行}\par

他又抓住四十个人，以新的方式殴打他们。他砍下棕榈树的鲜枝，上面仍带着刺，用它们猛烈抽打他们的脊背，以致其中一些人因刺断在肉里而长期接受外科治疗，另一些人因疼痛而死。他们把所有抓来的人和那个贞女，一同放逐到\Place{大绿洲}。那些死去者的尸体，他们起初不准交给亲友，而是随意隐藏，弃尸不葬，为要显得与这些残酷的行径无关。但他们的痴心妄想大大误导了他们。因为死者的亲友，既为宣认而欢欣，又为亲人的尸体悲伤，便将这作他们不虔与残忍的证据越发传扬出去。此外，他们立刻将下列主教逐出\Place{埃及}和\Place{吕彼亚}：\Person{阿摩尼乌斯}、\Person{缪伊乌斯}、\Person{加伊乌斯}、\Person{斐罗}、\Person{赫尔墨斯}、\Person{普莱尼乌斯}、\Person{普塞诺西里斯}、\Person{尼拉蒙}、\Person{阿加图斯}、\Person{阿纳甘福斯}、\Person{马尔谷}、\Person{阿摩尼乌斯}、另一个\Person{马尔谷}、\Person{德拉康提乌斯}、\Person{阿德尔菲乌斯}、\Person{阿特诺多鲁斯}，以及司铎\Person{希厄拉克斯}和\Person{狄奥斯库若}；他们是在如此残酷的待遇下被赶出的，以致一些人死在途中，其他人死在放逐地。他们还导致三十多位主教逃亡；他们的愿望是效法\Person{阿哈布}，若有可能，将真理彻底铲除。这些就是不虔之人所犯的暴行。\par

8. \Emph{若逃亡是错的，则迫害更恶。}\par

然而，尽管他们做了这一切，却对自己已向我策划的恶行毫不羞愧，反倒因为我得以逃脱他们杀人的毒手而继续控告我。岂料，他们倒为自己未能有效地将我除掉而痛苦哀叹；于是他们假装指责我胆怯，却未察觉，他们如此抱怨我，反倒将罪责转归自身。因为逃避若为坏事，迫害则更为恶劣；一方藏匿是为逃避死亡，另一方迫害则蓄意杀人；圣经上明载我们应逃避；但那图谋毁灭的人违犯法律，甚至正是他自己导致了对方的逃避。若他们因我逃避而指责我，就让他们对自己的迫害更感羞愧吧。让他们停止设谋，那些逃避的人便会立即停止逃避。然而，他们非但不停止其邪恶，反用尽一切手段要捉拿我，却未明白，受迫害者的逃避正是控告迫害者的有力论据。因为没有人会逃避温和仁慈之人，只会逃避残酷恶毒之辈。\BibleQuote{凡受窘迫的，和欠债的}\BibleRef{撒上 22:2}都逃离\Person{撒乌耳}，投奔\Person{达味}。而这些人之所以想剪除那些藏匿者，正是为了不让自己的恶行留下任何证据。但在这一点上，他们的心灵似乎因惯常的谬误而昏暗。因为敌人的逃避愈是广为人知，他们的奸计所带来的灭亡或放逐便愈昭然若揭；因此，他们或直接杀害他们，其死亡之声必将更响地传扬开来控告他们，或迫使他们流亡，则不过是到处竖立起自己邪恶的纪念碑。\par

9. \Emph{控告显示了控告者的心意。}\par

倘若他们心智正常，便应看出自己正处于此等困境，且正自相矛盾。然而，既然他们已丧失一切判断力，便仍然被驱策着去迫害、去毁灭，却从未察觉自己的不敬虔。他们甚至可能胆敢指责\OriginalTerm{天主眷顾}{Providence}本身（因为他们的僭越无所不及），怪它不将他们所觊觎的人交予他们；然而根据我们的救主的话，确凿无疑的是，没有天父的许可，连一只麻雀也不会落入网罗。\BibleRef{玛 10:29} 但当这些受诅咒者一旦捉住某人，他们便不仅忘记其他一切，甚至连自己都忘记了；他们狂妄地扬起额头，既不承认时代与时机，也不尊重他们所害之人的人性。就像\Place{巴比伦}的暴君，他们更加疯狂地攻击；对谁都不怜悯，反如经上所记，\BibleQuote{将重轭加在老人身上}，且\BibleQuote{增添受伤者的痛苦}。倘若他们不曾如此行事，未曾将那些为我辩护、反驳其诽谤的人放逐，他们的主张也许对某些人还显得颇为可信。但既然他们密谋陷害如此多德高望重的主教，既不放过那伟大的宣认者\Person{何西乌斯}，也不放过\Place{罗马}的主教，以及来自\Place{西班牙}、\Place{高卢}、\Place{埃及}、\Place{利比亚}和其他地方的众多人士，反倒对所有在维护我的事上以任何方式反对他们的人施以如此残忍的暴行；难道他们的图谋不是针对我甚于针对其他任何人，且他们那可怜的目的就是要像毁灭其他人那样毁灭我吗？为此，他们警惕地守候时机，眼见他们所不愿其存活的人安然无恙，便自觉受了伤害。\par

10. \Emph{他们真正的不满并非\Person{亚大纳削}怯懦，而是他自由了。}\par

如此，谁还认不出他们的诡诈呢？他们并非出于对德行的考虑而指责我怯懦，而是渴血成性，以这些卑劣伎俩为网罗，妄图借此捕捉他们所图谋毁灭之人。这岂非显而易见吗？他们的行为已显明他们的本性，证明他们具有比野兽更野蛮、比巴比伦人更残酷的性情。但尽管这一切已足够清晰地证明他们的罪状，他们却仍仿同\BibleQuote{父亲魔鬼}\BibleRef{若 8:44}的样式，用婉言欺瞒，假装指责我怯懦，而他们自己却比野兔更为怯懦；所以，让我们思考圣经对于这类情况是如何记载的。因为他们如此贬损圣人们的德行，便会显出他们不仅是在攻击我，更是在攻击圣经。\par

若他们指责人躲避那些图谋杀害他们的人，并控告逃离迫害者的人，那么当他们看到\Person{雅各伯}逃离他的哥哥\Person{厄撒乌}，\Person{梅瑟}因畏惧\Person{法郎}而退避到\Place{米德杨}时，他们又将如何呢？在这一切空谈之后，当他们看到\Person{达味}因\Person{撒乌耳}派人来杀他而从家中逃走，藏在洞穴里，并改变容貌，直到离开\Person{阿彼默肋客}而逃脱他的谋害时，他们又会作何辩解呢？当他们看到伟大的\Person{厄里亚}在呼求天主并复活死人之后，却因畏惧\Person{阿哈布}而躲藏，并躲避\Person{依则贝耳}的威胁时，这些随口妄议的人又将说些什么呢？那时，先知的门徒们被人搜捕，也因\Person{敖巴狄亚}的协助而藏匿，隐伏在洞穴中。\par

11. \Emph{圣经中圣人们为辩护逃避而立的榜样。}\par

也许他们没有读过这些史事，因为年代久远；可是他们难道不记得\WorkTitle{福音}中的记载吗？门徒们也曾因惧怕犹太人而退避躲藏；圣保禄在\Place{大马士革}被总督搜捕时，从城墙被装在篮子里缒下，这样逃过他的手。既然\WorkTitle{圣经}如此记述了圣人们的事迹，他们还有什么借口为自己的恶行辩解呢？指责圣人们胆怯乃是疯狂，控告他们行事违背天主的旨意，则表明自己完全不懂圣经。因为在法律中有过命令，吩咐设立避难城，好使那些被追捕处死的人至少能有自救之法\BibleRef{出 21:13}。而当那位对梅瑟说话者，即圣父的\OriginalTerm{圣言}{Logos}，在末世显现时，祂也赐下了这条诫命，说：\BibleQuote{当人们在这城迫害你们时，你们就逃往另一城去}；随后不久又说：\BibleQuote{几时你们看见达尼尔先知所说的\InnerQuote{招致荒凉的可憎之物}立于圣地，——读者应明白：——那时在犹太的，该逃往山中；在屋顶上的，不要下来从屋里取什么东西；在田野里的，也不要回去取他的外衣}。圣人们明白这些道理，便相应地规范了自己的行为。因为我们的主现今所命令的，也正是祂在降生成人之前藉着圣人们所说的：而这就是赐予人类以引领他们走向成全的规则——凡天主所命令的，就遵行。\par

12. \Emph{主是适时躲避的典范。}\par

因此，\OriginalTerm{圣言}{Logos}本身为了我们的缘故降生成人，也屈尊象我们一样在被人搜捕时躲藏起来；并在受迫害时逃走，以避开敌人的计谋。因为祂理当借着饥饿、口渴和受苦，也借着躲藏和逃走，显示出祂确实取了我们的肉身，成为了人。这样，从一开始，祂刚成为人，还是个小小婴孩时，就亲自通过祂的天使命令若瑟说：\BibleQuote{起来，带着婴孩和祂的母亲，逃往埃及，因为黑落德即将寻找这婴孩，要把祂杀掉}\BibleRef{玛 2:13}。当黑落德死后，我们又发现祂因黑落德的儿子阿尔赫劳而退往纳匝肋。后来，当祂显示自己是天主，并治好枯手的人时，法利塞人出去，商议怎样陷害祂，要将祂除掉；耶稣知道了，就从那里躲开了\BibleRef{玛 12:15}。同样，当祂使拉匝禄从死者中复活后，\WorkTitle{圣经}上说：\BibleQuote{从那天起，他们就议决要杀害祂。因此，耶稣不再公开地在犹太人中往来，而是从那里往靠近荒野的地方去了}\BibleRef{若 11:53}。还有，当我们的救主说：\BibleQuote{在亚巴郎出现以前，我就是}，\BibleQuote{犹太人便拿起石头来要砸祂；耶稣却隐没了，从圣殿里出去了}\BibleRef{若 8:58}。\BibleQuote{祂却由他们中间过去走了}\BibleRef{路 4:30}。\par

13. \Emph{我们的主的榜样。}\par

当他们看到这些事，或者甚至只是听到——因为他们没有看到——的时候，难道他们不会如经上所载，想要成为\BibleQuote{烈火的燃料}\BibleRef{依 9:5}吗？因为他们的计谋和言论都与主所行所教的背道而驰。在若翰蒙难，他的门徒埋葬了他的遗体后，\BibleQuote{耶稣听到了，就上船离开那里，独自到荒野地方去了}\BibleRef{玛 14:13}。主这样行了，也这样教导了。但愿这些人如今为自己的行为感到羞愧，把他们的狂妄限定在人的身上，而不要发展到如此极端疯狂的境地，竟至于指责我们的救主胆怯！因为他们现在口出亵渎的正是祂。但没有人会忍受这样的疯狂；相反，人们将看出他们不明白\WorkTitle{福音}。圣史们所记述的促使救主退避和逃走的理由，必定是一个合理而公正的理由；因此，我们也应该将此视为所有圣人的行事之理。（因为凡论到我们的救主在其人性方面的记载，都应该被视为适用于整个人类；因为祂取了我们的身体，并在自己身上展现了人性的软弱。）关于这个理由，若望这样写道：\BibleQuote{他们想要捉拿祂；但没有人下手，因为祂的时辰还没有到}\BibleRef{若 7:30}。在时辰来到之前，祂亲自对祂的母亲说：\BibleQuote{我的时辰还没有到}\BibleRef{若 2:4}；对那些被称为祂兄弟的人说：\BibleQuote{我的时候还没有到}\BibleRef{若 7:6}。后来，当祂的时辰来临时，祂对门徒们说：\BibleQuote{现在你们仍然睡觉安歇吧！看，时候到了，人子就要被交付在罪人手中}\BibleRef{玛 26:45}。\par

14. \Emph{为众人都有一个时辰和时期。}\par

就祂是天主及圣父的\OriginalTerm{圣言}{Word}而言，祂没有时间；因为祂自己是时间的创造者。但祂成为人，以这种方式说话，表明有一个时间是为每个人预定的；而且那不是偶然的，如同一些\OriginalTerm{外邦人}{Gentiles}在他们的神话中所想象的，而是那创造者按照圣父的旨意，为每一个人所指定的时间。这记载在圣经上，且为众人所明见。因为虽然每个人被指定的时间长度及其如何指定，对所有人都是隐秘未知的；然而每人都知道，如同有春天、夏天、秋天和冬天的时期一样，如经上所载\BibleRef{训 3:2}，有死时，也有生时。所以，诺厄时代的那一代人的时日被缩短，他们的年岁被减少，因为万有的时期近了。但希则克雅却又增添了十五年。正如天主应许那些真正事奉祂的人，\BibleQuote{我要满全你的寿数}\BibleRef{出 23:26}，亚巴郎\Quote{寿高年老}而死，达味曾恳求天主说：\Quote{求你不要使我中年逝世}。约伯的一位朋友厄里法次确知这真理，便说：\Quote{你必寿高年迈才归坟墓，好像禾捆准时收成}。撒罗满证实他的话，说：\Quote{不义者的灵魂夭折}。因此他在\WorkTitle{训道篇}中劝告说：\BibleQuote{不要过于邪恶，也不要成为愚顽：你为什么要在非时死去？}\BibleRef{训 7:17}\par

15. \Emph{主的时候与时辰}\par

既然这些事记载于圣经，情况就很清楚：圣人们知道每个人的时间是被度量好的，但无人知道那时间的终结，这由达味的话清楚暗示：\BibleQuote{求你使我明白我日子的短促}。他所不知道的，便切望得知。同样，那个富人，当他以为自己还有很长时间可活时，却听到这句话：\BibleQuote{糊涂人哪！今夜就要索回你的灵魂，你所准备的，将归谁呢？}\BibleRef{路 12:20}。\OriginalTerm{训道者}{Preacher}充满信心地在圣神内说：\BibleQuote{人也不知道自己的时期}\BibleRef{训 9:12}。因此\OriginalTerm{圣祖}{Patriarch}依撒格对他的儿子厄撒乌说：\BibleQuote{你看，我已年老，不知道我那天死的日子}\BibleRef{创 27:2}。所以，我们的主虽是天主和圣父的\OriginalTerm{圣言}{Word}，既知道祂为众人所量度的时间，也知道祂自己为祂的肉身所指定的受苦时辰；然而，祂为了我们成为人，在时辰未到之前，当祂被人搜寻时，便如我们一样隐藏自己；当祂被迫害时，便逃避了；祂避开敌人的计谋，过去了，并且\BibleQuote{从他们中间走过去}\BibleRef{路 4:30}。但当祂将祂自己所指定的那个时辰带来时，在那个时辰祂愿意为众人于肉身内受苦，祂便向圣父宣告说：\BibleQuote{父啊！时辰来到了，光荣你的圣子吧！}\BibleRef{若 17:1}。然后，祂不再躲避那些寻找祂的人，而是站着乐意被他们捉拿；因为圣经上说，祂对那些来到祂面前的人说：\BibleQuote{你们找谁？}\BibleRef{若 18:4}当他们回答：\BibleQuote{纳匝肋人耶稣}时，祂对他们说：\BibleQuote{我就是}。祂不止一次这样做了；于是他们立刻将祂带到了比拉多那里。在时辰未到之前，祂既不许自己被捉拿，时辰一到，祂也不隐藏；而是将自己交与那些图谋祂的人，好向众人显示：人的生死取决于天主的判决；若没有在天的父，人头上的一根头发也不能变白或变黑，一只麻雀也决不会落入罗网。\par

6. \Emph{圣人们效法主的榜样}\par

因此，我们的主，如我先前所说，这样为众人奉献了自己；圣人们从他们的救主领受了这榜样（因为所有人，在祂来临之前，不，无论何时，都在祂的教导之下），在与迫害者的斗争中，当被人搜寻时，他们合理地采取逃避和隐藏的做法。并且作为人，他们不知道\OriginalTerm{天主眷顾}{Providence}为他们所指定的时间终结如何，便不愿意立刻将自己交在那些图谋他们的人手中。但另一方面，他们知道经上这样记载：人的\Quote{分}在于天主手中；\BibleQuote{主使人死，也使人生}\BibleRef{撒上 2:6}；因此他们宁愿忍受到底，如宗徒所说：\BibleQuote{披着绵羊山羊的毛皮，穷乏、患难、受磨折，漂流于旷野，藏在山洞和地穴中}\BibleRef{希 11:37}；直到或者那指定的死期来临，或者那指定他们时间的天主向他们发言，阻止了他们敌人的计谋，或者按照祂所认为好的，将被迫害者交与迫害者。关于众人，我们可从达味身上清楚学到这一点：因为当约阿布怂恿他杀死撒乌耳时，他说：\BibleQuote{指着永生的主起誓：主必击杀他；或者他的死期到了；或者他下到战场，死于敌人之手。主绝不许我伸手加害主的受傅者。}\par

17. \Emph{逃避有时，留下有时}\par

他们在逃亡时若遇到追捕者，并非无缘无故；而是当圣神向他们发言时，他们便如义人一般前去迎见仇敌；藉此他们也显明了对天主的服从与热忱。\Person{厄里亚}的作为便是如此，他奉圣神之命去见\Person{阿哈布}\BibleRef{列上 21:18}；先知\Person{米加雅}去见同一个\Person{阿哈布}时亦然；还有那位在\Place{撒玛黎雅}祭台旁呼喊并责备\Person{勒哈贝罕}的先知；以及\Person{保禄}上告于\Person{凯撒}时。他们逃亡绝非出于怯懦：绝无此事。他们所承受的逃亡，实是一场与死亡的角力与争战。因为他们以明智的谨慎提防这两件事：要么无缘无故地献上自己（因为那无异于杀害自己，成为死亡的罪人，并违犯主的话：\BibleQuote{天主所结合的，人不可拆散}\BibleRef{玛 19:6}），要么甘愿遭受疏忽懈怠的指责，仿佛他们对逃亡中遭遇的苦难无动于衷，而这些苦难带来的痛苦比死亡更可怕、更惨重。因为死去的人不再受苦；但逃亡者日日等候敌人的攻击，反视死亡为轻。因此，那些在逃亡中结束尘世旅程的人，并非不光彩地丧亡，而是与别人一样获得了\OriginalTerm{殉道}{martyrdom}的荣耀。正因如此，\Person{约伯}被视为大有\OriginalTerm{勇德}{fortitude}的人，因为他忍受了如此多而沉重的苦难而活着，倘若他生命终结，便不会知晓这些苦难了。所以，真福的教父们也是如此安排自己的行为；他们逃避迫害者时毫不怯懦，反而将自己关闭在狭小阴暗之处，过着艰苦生活，以此彰显灵魂的勇德。然而，当死亡时刻来临时，他们并不渴望逃避；因为他们所关心的，既不是在它来临时刻意退缩，也不是抢先判决\OriginalTerm{天主眷顾}{Providence}所定的时刻，更不是反抗天主的安排——他们深知自己为此而蒙保全；以免行事急躁，成为自身恐惧的缘由：因为经上记载：\Quote{说话急躁的人，必自取惊吓。}\par

\Emph{逃亡的圣人们绝非懦夫。}\par

确实，无人能怀疑他们备有充足的勇德。因为圣祖\Person{雅各伯}先前曾逃避\Person{厄撒乌}，但当死亡来临时，他毫不惧怕，反而在那时依照各人的功过分赐祝福给圣祖们。伟大的\Person{梅瑟}，先前曾躲避\Person{法郎}，因惧怕他而退避到\Place{米德杨}；当他领受诫命\Quote{返回\Place{埃及}}时，毫不畏惧地遵行了。其后，当他受命前往\Place{阿巴陵山}\BibleRef{申 32:49}并死在那里时，他并不因怯懦而迟延，反而欣然前往。\Person{达味}先前曾逃避\Person{撒乌耳}，但在战争中为保卫人民，毫不畏死；当他面临死亡或逃亡的选择时，本可逃亡存活，却明智地选择了死亡。伟大的\Person{厄里亚}，早些时候曾躲避\Person{依则贝耳}，当他奉圣神之命去见\Person{阿哈布}并责备\Person{阿哈齐雅}时，毫不怯懦。\Person{伯多禄}曾因惧怕犹太人而躲藏，宗徒\Person{保禄}曾从篮子里缒下而逃亡，但当他们被告知\Quote{你们要在\Place{罗马}为我作证}时，他们毫不迟延旅程；相反，他们欢喜出发；一位仿佛急于会晤朋友，欣然接受死亡；另一位在死亡时刻来临时毫不退缩，反而夸耀说：\BibleQuote{因为我已被奠祭，我离世的时期已近了}\BibleRef{弟后 4:6}。\par

\Emph{圣人们在逃亡中勇敢无畏，又蒙天主恩宠。}\par

这些事既证明他们先前的逃亡并非出于怯懦；又表明他们后来的行为也非同寻常：它们大声宣告，他们拥有极高的勇德。因为他们既不是由于懒惰的畏缩而退避，相反，他们在这些时候比另一些时候更严格地修炼；他们也没有因为逃亡而被定罪，或被像如今热衷于指责他人的人指控为怯懦。不，他们因我主的这句话而蒙福：\BibleQuote{为义而受迫害的人是有福的}\BibleRef{玛 5:10}。这些苦难也并非对他们没有益处；因为就如智慧所说，天主试验他们，如同在炉中炼金，发现他们堪当属于自己\BibleRef{智 3:57}。然后，他们得脱迫害者之手，脱离敌人的计谋，被保全至终，好能教导百姓，因而越发闪耀如火花；所以，他们的逃亡与脱免追捕者的愤怒，乃按主的\OriginalTerm{安排}{dispensation}。这样，他们在天主眼中成为可爱的，并为他们的勇德领受了最荣耀的证言。\par

\Emph{同上继续。}\par

例如，圣祖雅各伯在逃亡时蒙恩获享许多甚至属神的异象，他自己保持安静，有主站在他一边，斥责拉班，挫败厄撒乌的计谋；后来他成为犹大的父亲，按血统来说，主正是从犹大而出；他又向众圣祖分赐了祝福。当天主所爱的梅瑟流亡时，就在那时他看见了大异象，他从迫害者手中被保全，被派遣作为先知进入埃及，成为那些大奇迹和法律的执行者，他带领那伟大的民族穿越旷野。达味受迫害时，写下了圣咏：\BibleQuote{我的心灵涌溢出美辞}；和\BibleQuote{我们的天主必要显现，决不缄默无声}。他再次更大胆地说：\BibleQuote{我的眼目看见了仇敌的报应}；又说：\BibleQuote{我全心仰赖天主，决不害怕世人能对我如何}。当他逃避撒乌耳的面，躲到\Quote{洞窟}时，他说：\BibleQuote{他从天上施救，拯救了我，使那企图践踏我的人蒙羞受辱。天主施予他的慈爱与忠诚，从狮子中间救出了我的灵魂}。这样，他也按天主的安排得蒙拯救，后来成为君王，并领受了应许，即我们的主将出自他的后裔。伟大的厄里亚退避到\Place{加尔默耳}山时，呼求天主，一举消灭了四百多个巴耳的先知；当有人派遣两个五十夫长及他们所带的百名士兵去捉拿他时，他说：\BibleQuote{愿火从天降下} \BibleRef{列下 1:10}，就这样斥责了他们。他也蒙了保全，因而得以代己傅油厄里叟，成为众先知弟子的纪律楷模。真福保禄在写下了这些话：\BibleQuote{我遭受了何等的迫害，但主从这一切当中救出了我，也必要救我} \BibleRef{弟后 3:11}；之后，更可坦然地说：\BibleQuote{但在这一切事上，我们已大获全胜，因为没有什么能叫我们与基督的爱隔绝} \BibleRef{罗 8:35}。因为正是在那时，他被提到第三层天，并被接入乐园，在那里听到了\BibleQuote{不可言传的话语，是人不能说出} \BibleRef{格后 12:4}。而正是为了这目的，他那时被保全，好能\BibleQuote{从耶路撒冷直到依里黎苛，将福音传得全备} \BibleRef{罗 15:19}。\par

21. \Emph{圣人为了我们而逃亡。}\par

因此，圣人们的逃亡既不应受指责，也并非无益。若他们不躲避迫害者，主如何出自达味的后裔呢？又有谁来宣讲真理之言的喜讯呢？迫害者搜捕圣人，正是为了使无人施教，正如犹太人指控宗徒们那样；但正因如此，他们忍受了一切，好使福音得以宣讲。因此，请看，他们如此与仇敌交锋，并非虚度逃亡的时光，也非在受迫害时忘记他人的益处：而是作为美善言语的仆役，他们不吝啬将之传于众人；以致即使在逃亡中，他们仍宣讲福音，警告那些谋害他们之人的邪恶，并以劝勉坚固信徒。因此，真福保禄亲身经历后，预先宣告说：\BibleQuote{凡愿意在基督耶稣内热心生活的人，都必遭受迫害} \BibleRef{弟后 3:12}。于是他立即鼓励那些为试炼而逃亡的人，说：\BibleQuote{让我们以坚忍，奔跑那摆在我们面前的赛程} \BibleRef{希 12:1}；因为虽然苦难不断，\BibleQuote{但是磨难生忍耐，忍耐生老练，老练生望德，望德不叫人蒙羞} \BibleRef{罗 5:4}。先知依撒意亚预期类似的磨难时，劝勉高呼说：\BibleQuote{我的百姓，去！进入你的内室，关上门，隐藏片刻，等待盛怒过去} \BibleRef{依 26:20}。同样，训道者知道针对义人的阴谋，说：\BibleQuote{你若看见穷人受压迫，并看见一省内审判与公义横遭歪曲，不要对此事惊奇；因为高处之上有更高者在鉴察，而他们之上还有更高者；况且，大地自有其利益}。他以自己的祖先达味为榜样，达味自己经历了迫害之苦，并用这些话支持受苦的人：\BibleQuote{鼓起勇气，你们所有寄望于主的人，他必坚固你们的心}；对那些这样坚忍的人，不是人，而是主自己（他说），\BibleQuote{必帮助他们，拯救他们，因为他们信靠他}；因为我也\BibleQuote{曾耐心等候主，他就垂听了我，俯听了我的呼求；他从祸坑与污泥中把我拉上来}。由此显明，圣人们的逃亡对人民何等有益且结出善果，不论\OriginalTerm{亚略派}{Arians}如何另作他想。\par

22. \Emph{同一主题的结语。}\par

因此，如我之前所说，圣人们在天主的眷顾下，在逃亡中得以丰盛保全，如同为有需要者而设的医生。而且普遍而言，对所有的人，甚至对我们，都给出了这条准则：遭受迫害时当逃亡，被搜寻时当躲藏，不可冒失试探主，而要等候，如我上面说的，直到预定的死亡时刻来到，或者审判者按他自己视为好的，对他们有所定断：人应做好准备，一旦时机召唤，或被拿获，就为真理奋斗到底，直至死亡。这一准则，真福殉道者在历次迫害中都遵守了。遭受迫害时，他们逃亡；隐藏时，他们显露刚毅；被发觉时，他们甘心殉道。如果有的人自行前来，将自己交给迫害者，他们这样做并非没有理由；因为那时他们立即殉道，从而向众人显明，这份热忱，以及将自己献给仇敌的行动，都是出于圣神。\par

23. \Emph{迫害来自魔鬼。}\par

既然我们\OriginalTerm{救主}{Saviour}的命令如此，\OriginalTerm{圣人}{Saints}们的行事亦是如此，就让那些无法用与其品行相称的名称称呼的人——我说，让他们告诉我们，他们是从谁那里学会迫害的？他们不能说是从圣人那里学的。不，而是从\OriginalTerm{魔鬼}{Devil}那里（这是他们唯一剩下的答案）——是从那位说\BibleQuote{我要追赶，我要追上}\BibleRef{出 15:9}的魔鬼那里学的。主命令我们逃遁，圣人们便逃遁；但迫害却是魔鬼的诡计，是他想要对所有人大施其技的计谋。那么，让他们说说，我们应当顺服哪一个：是主的话语，还是他们的臆造？我们应当效法谁的品行：是圣人们的，还是这些人的榜样？但他们似乎无法断定此问题（因为，正如\Person{依撒意亚}所说，他们的心思和良心已被蒙蔽，他们以为\BibleQuote{以苦为甜，以光为暗}\BibleRef{依 5:20}）。所以，让我们\OriginalTerm{基督徒}{Christians}中有人站出来，谴责他们，高声呼喊：\Quote{信赖上主，胜过听从这些人的愚言；因为主的话有永生的话\BibleRef{若 6:68}，而这些人所说的却充满罪孽与流血。}\par

\Emph{\Person{叙利亚努斯}的闯入}\par

这足以制止这些不虔敬者的疯狂，并证明他们除了一心争竞、口出辱骂和凌辱之外，别无他求。然而，既然他们胆敢对抗\Person{基督}，而今又如此殷勤，就让他们去询问自己的朋友，学习我撤离的方式吧。因为\OriginalTerm{亚略派}{Arians}混迹于士兵之中，为要激怒他们来攻击我，并因他们不识我的面貌，好将我指认给他们。虽然他们全无怜悯之情，但听见这些情形时，必将羞愧得哑口无言。那时已是夜间，有些人在为次日的共融而守夜祈祷，\Person{叙利亚努斯}将军突然带领超过五千名士兵，全副武装，手持出鞘的刀剑、弓、枪和棒，正如我上面所述。他用这些士兵包围了教堂，将士兵部署在近处，好使无人能够离开教堂从他们面前经过。我认为在这等骚乱中撇下百姓，而不为他们冒险，实为不近情理；因此我坐在我的座位上，吩咐\OriginalTerm{执事}{Deacon}颂读\OriginalTerm{圣咏}{Psalm}，并让百姓答唱：\Quote{因为祂的仁慈永远常存}，然后全体退去，各自回家。然而，将军此时强行进入，士兵包围了\OriginalTerm{圣所}{sanctuary}，为要捉拿我们，\OriginalTerm{圣职人员}{Clergy}和仍留在那里的\OriginalTerm{平信徒}{laity}就呼喊，要求我们也退去。但我拒绝了，声明除非他们全部退去，否则我绝不离开。于是我便站起来，吩咐祈祷，然后请求他们，都要在我之前离去，并说宁可我的安全受到威胁，也胜过他们中任何人受到伤害。如此，当大部分人都已离开，其余的人也正在跟上时，与我们在一起的\OriginalTerm{隐修士}{monks}们和某些圣职人员上前，将我们拖走。就这样（\OriginalTerm{真理}{Truth}是我的见证），当一些士兵守在圣所周围，另一些士兵在教堂内到处巡查时，我们却在主的引导下，靠着祂的庇护，不知不觉地经过，平安撤离，大力光荣天主，因为我们没有背弃百姓，反而先遣散了他们，而后又得以保全自己，逃脱了搜捕我们之人的手。\par

\Emph{\Person{亚大纳削}的奇妙逃脱}\par

如今\OriginalTerm{天主眷顾}{Providence}既以如此奇特的方式拯救了我们，谁还能公正地指责我，因我们没有将自己交在那搜捕我们之人的手中，没有回去主动出现于他们面前呢？这分明是对主忘恩负义，违背祂的命令，与圣人们的做法背道而驰。凡在此事上指责我的人，无疑也必指责伟大的宗徒\Person{伯多禄}，因为他虽被关押并有兵士看守，却跟随了召叫他的天使，而且在他逃出监牢、平安脱险之后，并没有返回去自首，尽管他听到了\Person{黑落德}所行的事。在狂怒中的亚略派，就让他去指责宗徒\Person{保禄}，因为他从城墙上被缒下脱险之后，并没有改变心意回去自首；或指责\Person{梅瑟}，因为他没有从\Place{米德杨}返回\Place{埃及}，以便让搜捕他的人擒获；或指责\Person{达味}，因为他藏匿在洞中时，并没有向\Person{撒乌耳}显身。同样，先知们的弟子们也留在洞穴里，也没有向\Person{阿哈布}投降。这样做便是违背诫命，因为经上说：\BibleQuote{你不可试探上主你的天主。}\par

\Emph{他按圣人的榜样行事。控告者的品性}\par

我谨慎地避免这种冒犯，并从这些榜样中得到教训，便如此安排了我的行为；我绝不低估主所赐予我的恩惠和帮助，无论这些人在疯狂中如何对我们咬牙切齿。因为我们的退避方式既已知是如我们所描述的那样，我认为凡具有健全判断力的人，都不会因此归咎于我们：因为按照\OriginalTerm{圣经}{Scripture}，\OriginalTerm{圣人}{Saints}们已为我们留下了这样的榜样，供我们效法。然而，看来这些人没有不行的暴行，他们也不会放过任何能显露自己邪恶与残忍的事。的确，他们的生活只是与其精神和教义的愚妄相称；因为无论多么可憎的罪恶，他们无不厚颜无耻地犯下。\Person{雷昂提乌斯}因与一位名叫\Person{欧斯托利乌}的年轻女子过从甚密而受责，并被禁止与她同居，他竟为她而自阉，以便能自由地与她交往。然而他并未因此洗脱嫌疑，反而因此被罢免了\OriginalTerm{司铎}{Presbyter}圣职。[尽管\OriginalTerm{异端者}{heretic}\Person{君士坦提乌斯}以暴力使他被称为\OriginalTerm{主教}{Bishop}。]\Person{纳尔基索斯}除了被指控犯有许多其他罪行外，还曾三次被不同公会议罢免；现在他仍混迹于他们中间，是个极其邪恶的人。而\Person{乔治}曾是一位司铎，因邪恶被罢免，尽管他自封为主教，却仍在\Place{萨迪卡}大公会议上被第二次罢免。此外，他放荡的生活声名狼藉，因为甚至他自己的友人也谴责他，认为生存的终极目的与幸福在于犯下最可耻的罪行。\par

27. \Emph{结论}\par

如此，各人都在自己特有的恶习上彼此更胜一筹。但他们全都有一个共同的污点，即因着自己的异端，他们与基督为敌，不再被称为基督徒，而是亚略派。他们本应彼此控告所犯的罪，因为那些罪违背了基督的信仰；但他们却为了自己的缘故而遮掩这些罪。毫不奇怪，他们既具有这样的精神，又陷于如此罪恶之中，便迫害并追捕那些不随从他们同样不虔的异端的人；他们乐于毁灭这些人，若不能得逞便忧伤，正如我先前所说，当他们看到渴望其死的人还活着时，便觉得自己受了委屈。愿他们继续如此受委屈，甚至丧失伤害他人的能力，而他们所迫害的人能感谢主，并用第二十六篇\BibleBook{}{圣咏}的话说：\BibleQuote{上主是我的光明，我的救援，我还畏惧何人？上主是我生命稳固的保障，我还害怕何人？当恶人，就是我的仇敌和敌对我的人，前来攻击我，要吃我肉时，他们便绊倒跌下；}\BibleRef{咏 26:1-2} 又说第三十篇\BibleBook{}{圣咏}：\BibleQuote{祢救我的灵魂脱离危难，没有将我交于敌手；祢使我的脚立于广宽之处。}\BibleRef{咏 30:8} 在基督耶稣我们的主内，藉着祂，在圣神内，愿光荣和权能归于圣父，至于无穷之世。阿们。\par

\SourceRecord{Translated by M. Atkinson and Archibald Robertson.；Nicene and Post-Nicene Fathers, Second Series；Vol. 4.；Edited by Philip Schaff and Henry Wace.；Buffalo, NY: Christian Literature Publishing Co.,；1892.；<http://www.newadvent.org/fathers/2814.htm>.}
